\textbf{Expansao de Taylor de $\sin x$ em torno de $x_0=0$ (Maclaurin), com $n=4$:}

A formula de Taylor com resto de Lagrange e:
\[
f(x)=P_4(x)+R_5(x),
\]
onde
\[
P_4(x)=f(0)+f'(0)x+\frac{f''(0)}{2!}x^2+\frac{f^{(3)}(0)}{3!}x^3+\frac{f^{(4)}(0)}{4!}x^4,
\]
\[
R_5(x)=\frac{f^{(5)}(\xi)}{5!}x^5,\qquad \xi\in(0,x).
\]

\textbf{O que e $R_5(x)$, em palavras simples?}

\begin{itemize}
\item $P_4(x)$ e a parte que usamos para aproximar $\sin x$.
\item $R_5(x)$ e o que ficou de fora da aproximacao.
\item Portanto, $R_5(x)$ e exatamente o erro:
\[
R_5(x)=\sin x-P_4(x).
\]
\end{itemize}

Ou seja, se $R_5(x)$ for pequeno, a aproximacao por Taylor e boa.

Para $f(x)=\sin x$:
\[
f(0)=0,\quad f'(0)=1,\quad f''(0)=0,\quad f^{(3)}(0)=-1,\quad f^{(4)}(0)=0.
\]

Logo, o polinomio de Taylor de ordem 4 e:
\[
P_4(x)=x-\frac{x^3}{3!}=x-\frac{x^3}{6}.
\]

Assim, a aproximacao pedida fica:
\[
\sin x\approx x-\frac{x^3}{6},
\]
e o erro e exatamente o resto:
\[
R_5(x)=\frac{f^{(5)}(\xi)}{5!}x^5=\frac{\cos(\xi)}{120}x^5.
\]

Como $|\cos(\xi)|\le 1$, obtemos a cota:
\[
|R_5(x)|\le \frac{|x|^5}{120}.
\]

\textbf{Exemplo rapido:} para $x=0{,}5$,
\[
|R_5(0{,}5)|\le \frac{0{,}5^5}{120}=2{,}60\times10^{-4}.
\]
Isso quer dizer que o erro ao usar $\sin(0{,}5)\approx 0{,}5-0{,}5^3/6$ e, no maximo, da ordem de $10^{-4}$.

\textbf{Exemplos numericos da cota do erro:}
\begin{center}
\begin{tabular}{ll}
\hline
$x$ & $|R_5(x)|\le |x|^5/120$ \\
\hline
$0{,}1$ & $8{,}33\times 10^{-8}$ \\
$0{,}5$ & $2{,}60\times 10^{-4}$ \\
$1{,}0$ & $8{,}33\times 10^{-3}$ \\
$\pi/4$ & $2{,}49\times 10^{-3}$ \\
\hline
\end{tabular}
\end{center}

\textbf{Conclusao:} com $n=4$, o erro da aproximacao de $\sin x$ por $x-\frac{x^3}{6}$ e da ordem de $x^5$ e fica limitado por $|x|^5/120$, sendo muito pequeno para $|x|<1$.